\section*{Limitations}
\label{sec:limitations}

\paragraph{One training run per arm.}
Every effect size reported here is computed across the 96 patient personas \emph{within} a single
training run. There are no training-seed replicates, so run-to-run training variance is unmeasured,
and we cannot separate ``look-ahead helps'' from ``this look-ahead run went well.'' This is the
limitation a reader should weight most heavily. It is distinct from, and not bounded by, the
evaluation noise floor below.

\paragraph{Evaluation draws: replicated at the endpoint, single elsewhere.}
Therapist decoding is unseeded and each model state was evaluated on one 96-conversation sample.
Two noise floors bound this. At the untrained base, four independent draws of the identical policy
yield 54 same-policy contrasts with none reaching even uncorrected $p < .05$, but that is measured
where sessions are short and homogeneous. At the \emph{trained} endpoint we re-drew the winning
state (\S\ref{sec:reward}): 9 instruments $\times$ 2 graders of the same policy differ by at most
$|\dz|\,0.174$, nothing significant, and the $K$ contrast keeps its sign, significance, and effect
size. The remaining states are single draws, and a replicate bounds evaluation noise only, never
the training variance above.

\paragraph{Matched iterations are not matched cost.}
Every comparison is at matched iteration counts. An iteration is 96 conversations and $G{=}8$
scored candidates per prompt in both arms, so oracle scoring calls are approximately matched by
construction ($302{,}541$ for $\Kz$ vs.\ $289{,}983$ for $\Kf$ over ten iterations). But each
$\Kf$ reward requires a five-turn simulated continuation, about $393$k patient-simulator calls
over the run that the $\Kz$ arm never makes, and roughly doubles per-step wall-clock (median
$1.92\times$ over the settled iterations 3--10). More per iteration is built into the lever; we
disclose the asymmetry rather than equalise it, and nothing here compares the arms at matched cost.

\paragraph{$K \in \{0, 5\}$ only.}
There is no dose--response evidence between the two levels and nothing above 5, so nothing here
says the effect is monotone in $K$, that 5 is optimal, or that a shorter horizon would not capture
most of the benefit. An intermediate $K$ is the most valuable single follow-up.

\paragraph{The look-ahead reward is also a reward for continuing.}
A simulated session can end inside the look-ahead window, most often because the patient closes
it, and such a candidate gives the oracle a shorter transcript. Over all $121{,}088$ scored $\Kf$
candidates, $82\%$ of rollouts ran the full five turns (the early-ending share ranged from $12\%$
at iteration 10 to $30\%$ at iteration 6, and in $16\%$ of candidates the patient closed);
early-ending candidates scored at or below their group mean, by up to $0.09$ points depending on
how the rollout ended, and were the group's best less often than chance ($0.10$ vs.\ $0.125$;
Appendix Figure~\ref{fig:tails}). The lever was therefore applied as specified in the large
majority of cases, but it carries a mild pressure to keep the session open that a turn-level
reward does not, and the $\Kf$ policy's longer endpoint sessions (\S\ref{sec:behaviour}) are
consistent with it. We cannot separate this from the horizon effect proper with two arms.

\paragraph{Both policies grew into the response cap.}
Therapist responses are capped at 200 tokens, and mean turn length roughly tripled over training in
both arms, so many endpoint turns end mid-sentence (Appendix~\ref{app:example}). The cap is
identical across arms and the truncated turn is what the patient simulator and both graders saw,
so it does not bias the contrast, but it is part of what every score in this paper measures.

\paragraph{Every number is in-sample with respect to the patient distribution.}
The 96 personas are the exhaustive enumeration of the persona grid, and the self-training loop uses
those same 96 as both the training rollouts and the evaluation set at every iteration.
Generalisation to unseen patients is not merely unmeasured but unmeasurable without authoring new
personas.

\paragraph{The grader is decoupled; the generator is not.}
The held-out judge addresses the concern that the patient simulator and the training oracle are the
same model, but only for the grading half. Every conversation was generated against a
\texttt{gpt-4o-mini} patient. A fully decoupled replication would regenerate all conversations with
a patient from another model family.

\paragraph{No human validation.}
All eight instruments are administered by language models. The oracle is highly self-repeatable
(ICC 0.86--0.99 on the measured $K{=}0$ subset), but repeatability is not validity. No human MI
coder has rated any conversation in this experiment. The over-praise finding is the most robust
behavioural claim here precisely because a deterministic lexical marker corroborates it without a
model in the loop.

\paragraph{The reward is also a reported outcome.}
\QQ\ is both the optimisation target and a headline metric, so ``\QQ\ improved'' is partly
circular. We treat the held-out judge and the instruments outside the reward as the load-bearing
evidence, and \S\ref{sec:measurement} quantifies what happens to the trained-against grader's
discriminative range over a long run.

\paragraph{Instrument-specific reliability.}
The MITI-style coder is the least dependable instrument at the arm level, and we do not lean on
its differences alone. There is no reliability estimate for any individual MI-inconsistency
channel, including over-praise, which is the granularity \S\ref{sec:behaviour} argues at; that is
the main reason the section's conclusion rests on the deterministic marker.
